\setlength{\headheight}{23pt}
\usepackage[a4paper, left=2.5cm, right=2.5cm]{geometry}
\usepackage[utf8]{inputenc}
%\usepackage[T1]{fontenc}
\usepackage[ngerman]{babel}
\usepackage{grffile}
\usepackage{eurosym}
\usepackage{graphicx} % Allows including images 
\usepackage{amsmath, amsthm, amssymb}
\usepackage[separate-uncertainty]{siunitx}
\usepackage{ifthen}
\usepackage{fancyhdr}
\pagestyle{fancy}
\usepackage[justification=centering,labelfont=bf]{caption}
\usepackage[bottom]{footmisc}
\usepackage{todonotes}
\presetkeys{todonotes}{inline}{}
\usepackage{tikz}
\usepackage{gnuplot-lua-tikz}
\usepackage{circuitikz}
\usetikzlibrary{calc}
\usepackage{lastpage}
\usepackage{xspace}	
\usepackage[skip=10pt plus1pt, indent=10pt]{parskip}
\usepackage{makecell}
\usepackage{csvsimple}
\usepackage{siunitx}
\usepackage{svg}
\usepackage{sectsty}
\allsectionsfont{\sffamily}
\usepackage{csquotes}
\usepackage{multicol}
\usepackage{setspace}

\usepackage{enumitem}
\setitemize{itemsep=-5pt}
\setenumerate{itemsep=-5pt}
\setlength{\mathindent}{10pt}
\renewcommand{\baselinestretch}{1.5} 

\newcommand\tab[1][0.5cm]{\hspace*{#1}}

\usepackage[					   % Literaturverzeichnis mit "biber [Projektname ohne Endung]" aktualisieren
    backend=biber,
    style=ieee,
]{biblatex}                        % Bibliography management
\addbibresource{references.bib}    % Bibliography file

%%%%%%%%%%%%%%%%%%%%%%%%%%%%%%%%%%%%%%%%%%%%%%%%%%%%%%%%%%%%%%%%%%%%%%%%%%%%%%%%
%%%%%%% PACKAGAES %%%%%%%%%%%%%%%%%%%%%%%%%%%%%%%%%%%%%%%%%%%%%%%%%%%%%%%%%%%%%%
%%%%%%%%%%%%%%%%%%%%%%%%%%%%%%%%%%%%%%%%%%%%%%%%%%%%%%%%%%%%%%%%%%%%%%%%%%%%%%%%

%  Mathematik
\usepackage{amsmath}								%  Mathematikpaket
\allowdisplaybreaks									%  Formeln können auf mehrere Seiten verteilt werden
\usepackage{amssymb}								%  spezielle Mathmematiksymbole
\usepackage{siunitx}								%  korrektes Darstellen von Einheiten
\sisetup{locale = DE}

\usepackage{geometry}								%  Definition der Seitenraender
\usepackage[final]{pdfpages}							% zum einbinden von externen pdf-Seiten
\usepackage{graphicx}								%  zum Einbinden von Grafiken, spziell für dvi->ps Konvertierung
\setlength{\unitlength}{1mm}
\usepackage[hidelinks]{hyperref}	
\usepackage{placeins}								% FloatBarrier
\usepackage{upgreek}
\usepackage{tabularx}								% schöne tabellen
\usepackage{longtable}  								% lange tabellen
%\usepackage[osf,sc]{mathpazo}						% Schriftart	
\usepackage{booktabs}								% linien in der tabelle
\usepackage{nicefrac}								% brüche in zeilen

\usepackage{array}
\newcolumntype{Y}{>{\centering\arraybackslash}X}


% ---------------Mathematisches----------------------------->
%  fett dargestellte Zeichen im Mathematikmodus, Matrizen, Vektoren
\newcommand{\mv}[1]{\ensuremath{\boldsymbol{#1}}}
%  fett dargestellte unterstrichene Zeichen im Mathematikmodus
\newcommand{\kmv}[1]{\ensuremath{\underline{\boldsymbol{#1}}}}
%  Transponierte einer Matrix
\newcommand{\trans}[1]{\ensuremath{\boldsymbol{#1}^{\mathrm{T}}}}
%  Transponierte der konjugiert komplexen Matrix
\newcommand{\herm}[1]{\ensuremath{\boldsymbol{#1}^{\mathrm{H}}}}
%  Rang Matrix
\newcommand{\rg}[1]{\ensuremath{\mathrm{rg}\left({#1}\right)}}
%  Spur einer Matrix
\newcommand{\tr}[1]{\ensuremath{\mathrm{tr}\left\{ #1 \right\}}}
%  Realteil
\newcommand{\re}[1]{\mathrm{Re}\left\{  #1 \right\} }
%  Imaginärteil
\newcommand{\im}[1]{\mathrm{Im}\left\{  #1 \right\} }
%  Unterstreichung ... komplex
\newcommand{\K}[1]{\underline{#1}}
%  erzeugt das "entspricht Zeichen"
\newcommand{\corsp}{\ensuremath{\stackrel{\scriptscriptstyle \bigtriangleup}{=}}}
%  gewöhnliche Ableitung
\newcommand{\ddx}[1][x]{\frac{\mathrm{d}}{\mathrm{d} #1}}
%  partielle Ableitung
\newcommand{\pddx}[1][x]{\frac{\partial}{\partial #1}}
%  zweifache partielle Ableitung
\newcommand{\ppddx}[1][x]{\frac{\partial^2}{\partial #1 ^2}}
%  d-Operator
\newcommand{\dd}{\mathrm{d}}
%  Einheitsmatrix							
\newcommand{\ident}{\mv{I}}	
%  const. im Mathematikmodus							
\newcommand{\const}{\mathrm{const.}}			
%  imaginäre Einheit
\newcommand{\ii}{j}		
%  imaginäre Einheit * omega										
\newcommand{\jw}{\ii \, \omega}						
% <--------------- End of Mathematisches ----------------------

% --------------- Abkürzungen ------------------------------>
\newcommand{\zb}{\mbox{z.\,B.}\xspace}		%  Erzeugt z.B. mit korrektem Zwischenraum
\newcommand{\ua}{\mbox{u.\,a.}\xspace}		%  Erzeugt u.a. mit korrektem Zwischenraum
\newcommand{\ie}{\mbox{d.\,h.}\xspace}		%  Erzeugt d.h. mit korrektem Zwischenraum
\newcommand{\og}{\mbox{o.\,g.}\xspace}		%  Erzeugt o.g. mit korrektem Zwischenraum
\newcommand{\oa}{\mbox{o.\,a.}\xspace}		%  Erzeugt o.a. mit korrektem Zwischenraum
% <--------------- End of Abkürzungen ----------------------

\lhead{} 
\chead{} 
\rhead{} 

\lfoot{} \cfoot{} \rfoot{Seite \thepage\ von \pageref{LastPage}} 
\renewcommand{\headrulewidth}{0.4pt}
\renewcommand{\footrulewidth}{0.4pt}
\addto\captionsngerman{\renewcommand{\figurename}{Abb.}}
\addto\captionsngerman{\renewcommand{\tablename}{Tab.}}
%\captionsetup{justification = raggedright}
